\documentclass[output=paper,colorlinks,citecolor=brown]{langscibook}
\ChapterDOI{10.5281/zenodo.18845594}
\author{Maria Petrou\orcid{}\affiliation{} and Hamza Boutemin\orcid{}\affiliation{} and Axel Fanego Palat\orcid{}\affiliation{}}
\title[Perceiving communicative behavior as discriminatory]{Perceiving communicative behavior as discriminatory: Tarifit and German speakers' experiences}
\abstract{
%This chapter explores how North African men in Germany discuss their experiences of perceived discrimination, focusing on the role of group boundaries and affiliations in shaping responses to verbal discrimination. The study analyzes a group conversation involving three North African and two German-background men, examining how linguistic differences, including the use of multiple languages, influence communication dynamics and intergroup relations. Drawing on Communication Accommodation Theory (CAT) and critical intercultural communication frameworks, the chapter looks at how speakers adjust their linguistic expressions and behaviors in multilingual settings.

%The findings indicate that group boundaries such as ethnicity, language, and social status impact participants’ perceptions of discrimination and their identity negotiations. Despite shared multilingual competence, participants navigate asymmetries in language proficiency, social prestige, and power dynamics, leading to misunderstandings and nonaccommodation, which can reinforce stereotypes. The study also highlights how stereotypes are contested, adapted, or subverted in everyday communication, with some derogatory terms being used as expressions of solidarity.

%The chapter emphasizes that nonaccommodation in intergroup communication often reflects deeper linguistic, social, and power imbalances, rather than simply misperceptions. It advocates for a dual approach combining CAT and critical discourse analysis\is{critical discourse analysis (CDA)} (CDA) to better understand intercultural and intergroup communication in the context of discrimination. Such an analysis provides a broader conceptual toolkit, offering richer insights into the dynamics of discrimination and the emotional impact of nonaccommodative behaviors.

This chapter explores how North African men in Germany discuss perceived discrimination, focusing on group boundaries and affiliations shaping responses to verbal discrimination. Analyzing a group conversation with three North African and two German-background men, the study examines how linguistic differences, including multilingual use, influence communication dynamics and intergroup relations. Drawing on Communication Accommodation Theory (CAT) and critical intercultural frameworks, it investigates linguistic adjustments and behaviors in multilingual settings.

Findings indicate that group boundaries (ethnicity, language, social status) impact discrimination perceptions and identity negotiations. Despite shared multilingual competence, participants navigate asymmetries in language proficiency, social prestige, and power dynamics, leading to misunderstandings and nonaccommodation that reinforce stereotypes. The study also highlights how stereotypes are contested, adapted, or subverted in everyday communication, with some derogatory terms used as expressions of solidarity.

The chapter emphasizes that nonaccommodation reflects deeper linguistic, social, and power imbalances rather than mere misperceptions. It advocates combining CAT and critical discourse analysis\is{critical discourse analysis (CDA)} to better understand intercultural communication in discrimination contexts. This dual approach provides a broader conceptual toolkit, offering richer insights into discrimination dynamics and the emotional impact of nonaccommodative behaviors.}

\IfFileExists{../localcommands.tex}{
   \addbibresource{../localbibliography.bib}
   \usepackage{tabularx,multicol}
\usepackage{url}
\urlstyle{same}
\usepackage{textcomp}
 
\usepackage{langsci-optional}
\usepackage{langsci-lgr}
\usepackage{langsci-gb4e}
\usepackage{langsci-branding}

% \usepackage{fontspec}
\babelprovide{georgian} % Georgian language
\babelprovide{armenian} % Armenian language
\babelfont[georgian,armenian]{rm}{DejaVu Sans}
\newfontfamily\georgianfont{DejaVu Sans}
\newfontfamily\armenianfont{DejaVu Sans}

\usepackage{dialogue}
\usepackage{attrib}
\usepackage{attrib} % speaker labels in Koca's chapter
\usepackage{phoenician} % alf, bet in Mengozzi's chapter

   

\makeatletter
\let\thetitle\@title
\let\theauthor\@author
\makeatother

\newcommand{\togglepaper}[1][0]{
   \bibliography{../localbibliography}
   \papernote{\scriptsize\normalfont
     \theauthor.
     \titleTemp.
In Saloumeh Gholami \& Ulrich Mehlem (eds.), \textit{Identity formation, language and migration dynamics: Minorities in the MENA region and the German diaspora}, Berlin: Language Science Press [preliminary page numbering]
   }
   \pagenumbering{roman}
   \setcounter{chapter}{#1}
   \addtocounter{chapter}{-1}
}

\newbool{bookcompile}
\booltrue{bookcompile}
\newcommand{\bookorchapter}[2]{\ifbool{bookcompile}{#1}{#2}}


\newcommand{\georgian}[1]{{\georgianfont #1}}
\newcommand{\armenian}[1]{{\armenianfont #1}}

\newcommand{\textarab}[1]{{\arabicfont \beginR #1} \endR}
\newcommand{\texthebrew}[1]{{\hebrewfont \beginR #1} \endR}
\newcommand{\textsyriac}[1]{{\syriacfont \beginR #1} \endR}


\newfontfamily\phoenicianfont{Tuffy Regular.ttf}

   %% hyphenation points for line breaks
%% Normally, automatic hyphenation in LaTeX is very good
%% If a word is mis-hyphenated, add it to this file
%%
%% add information to TeX file before \begin{document} with:
%% %% hyphenation points for line breaks
%% Normally, automatic hyphenation in LaTeX is very good
%% If a word is mis-hyphenated, add it to this file
%%
%% add information to TeX file before \begin{document} with:
%% %% hyphenation points for line breaks
%% Normally, automatic hyphenation in LaTeX is very good
%% If a word is mis-hyphenated, add it to this file
%%
%% add information to TeX file before \begin{document} with:
%% \include{localhyphenation}
\hyphenation{
    par-a-digm
    Ver-to-vec
    Go-go-lin
    sa-gen
    mo-misch
    ant-wor-te-te
    dix-wa-zim
    spre-chen
    B-Verf-GE
    Swe-dish
    re-mains
    schrei-ben
    di-pey-vin
    Chris-ten-sen
    Wes-sen-dorf
    Ja-va-khi-shvi-li
    ana-lysis
}

\hyphenation{
    par-a-digm
    Ver-to-vec
    Go-go-lin
    sa-gen
    mo-misch
    ant-wor-te-te
    dix-wa-zim
    spre-chen
    B-Verf-GE
    Swe-dish
    re-mains
    schrei-ben
    di-pey-vin
    Chris-ten-sen
    Wes-sen-dorf
    Ja-va-khi-shvi-li
    ana-lysis
}

\hyphenation{
    par-a-digm
    Ver-to-vec
    Go-go-lin
    sa-gen
    mo-misch
    ant-wor-te-te
    dix-wa-zim
    spre-chen
    B-Verf-GE
    Swe-dish
    re-mains
    schrei-ben
    di-pey-vin
    Chris-ten-sen
    Wes-sen-dorf
    Ja-va-khi-shvi-li
    ana-lysis
}

   \boolfalse{bookcompile}
   \togglepaper[2]%%chapternumber
}{}
 
\begin{document}
\maketitle

\section{Analyzing nonaccommodation critically}\label{CH1.1.ss.1}

Our chapter intends to shed light on how group boundaries and affiliations play out in how especially verbal discrimination is perceived, reacted to, and also ``put into practice''. We invite readers to reflect on notions of ingroup versus outgroup, their significance for what is said, and how it is talked about in a small group of five men of North African and \ili{German} backgrounds. North African men living in Germany talk about their experiences of being discriminated against. We draw on a careful analysis of a group conversation among these men, three originally from North Africa and two of \ili{German} background. Our theoretical interest revolves around the question of how speakers adjust to one another when more than one linguistic code is involved. Intercultural communication is inevitably more prone to misunderstandings and misperceptions. These threaten the functioning of mutual adjustments of linguistic expressions, styles, speech rate, word choice, and numerous other features of (linguistic) communication \citep{giles_bourhis_taylor1977}.
%\largerpage[2]

Communicating across language boundaries implies that at least one party might be placed at a disadvantage if the conversation takes place in the other party's dominant language. Often this goes hand in hand with asymmetries in prestige, social status and a sense of entitlement to positions of power. However, in our view, we should not frame our analysis simply in terms of ``lack of mastery'' of different, supposedly clearly delineated linguistic codes in the sense of distinct languages. We should rather keep an eye on the repertoires available in each case. In this respect, the approach differs from the ``classical'' work on foreign(er) talk (e.g., \citealt{zuengler1991}). We consider this to be meaningful and necessary in order to be able to capture precisely those cases in which shared multilingual competence is mutually assumed by the interlocutors, while they apply the various codes differently in order to occupy and claim discursive positions in intergroup communication.

We need to concern ourselves with the question of how groups are constituted and perceived and what their respective significance is for different participants in the discourse. Therefore, we contrast episodes from different contexts, all drawn from the recorded group conversation. In the main episode, a participant perceived inappropriate communicative behavior by an ingroup member in the Rif region (for further details, see \ref{3.8.ss.2}). Further insight is drawn from brief episodes concerning nonverbal behavior by members of the outgroup, in this case, in Germany. These latter incidents are small but frequent moments at which group boundaries and expectations to be discriminated against frame the perceptions of this study's participants: other people's behavior that is interpreted as a reaction to one's mere presence as a young male migrant. In these cases, it is not exactly clear what the relevant groups for those involved in the encounters are. Whenever groups are named entities, though, their existence has obviously gained sufficient currency to deserve a label -- ethnic, linguistic, cultural, socio-economic or other. Being a member of such a group invites inferences about oneself by others based on their general assumptions concerning that group. Stereotyping is a common and, to a considerable extent, necessary feature (\citealt{lakoff1987}) when humans construe conceptual categories such as ``Arabs'', ``refugees'', and others that we find mentioned in the conversation on which this chapter is based. At the same time, these categories are typically contested and not unequivocal -- contrary to their general perception. They are constantly appropriated, rejected, adapted or subverted by members of the respective groups. Their labels are employed for banter or as insults, but their offensive potential can be turned into the opposite, e.g., when derogatory words are used as terms of endearment under the right communicative conditions. 

Our study deals both with content-related topics as brought up by our interlocutors, as well as with linguistic norms and their application or manipulation by our interlocutors. Our methodology strives to link both aspects. We choose a specific conversation setting: a group interview on experiences of discrimination that we analyze in-depth relying on categories that are inductively identified in the material itself \citep{kuckartz2016,mayring2016}. We identify particular cases in our data that involve instances of non-optimal communication between the interlocutors. These include examples of experiences reported by the interviewees, but also concrete examples from the recorded conversation situation itself. We analyze these instances following a \isi{Communication Accommodation Theory} (CAT) framework \citep{giles2016}.

\begin{sloppypar}
We suggest the following hypothesis: In intergroup communication, nonaccommodation is not only associated with the interlocutors' incorrect perceptions and interpretations that lead to misunderstandings and unintentional deviations from the communicative maxim of accommodation. Multilingual\is{multilingual} situations are often characterized by asymmetries with regard to the interlocutors' linguistic knowledge of certain codes, with regard to the prestige associated with the codes, but also with regard to the evaluation of the content-meaning of certain topics.
\end{sloppypar}

\section{Participants in a conversation about discriminatory practices}\label{CH1.1.ss.2}

The experience of being discriminated against was the main topic of the group interview that formed the basis of this contribution. It involved five participants, three of North African background and two from Germany. In addition to the authors Boutemin and Fanego Palat, two young North African men, Hilal and Lamin,\footnote{%
For the participants, excluding the authors, we assigned pseudonyms to keep their identity anonymous.} and a \ili{German} student, Marcus, participated in the interview. Hilal and Lamin were the main protagonists of the conversation that we recorded and analyzed. Hilal and Lamin are from Nador and speak Tarifit\il{Berber!Tarifit} (``Rif \ili{Berber}''), an Amazigh\il{Berber!Amazigh|see{Berber!Tamazight}} variety used mainly in the Rif area in \isi{Morocco} (but with a substantial number of speakers across European countries due to migration). Lamin self-identifies as a \ili{Darija}-Tarifit\il{Berber!Tarifit} bilingual due to his upbringing in a mixed family (Lamin [00:04:23--00:05:00]). Hilal describes himself primarily as \ili{Berber}-speak\-ing,\footnote{%
When speaking \ili{German}, Hilal commonly (although not exclusively) refers to his first language as ``\ili{Berber}'', sometimes as ``Tamazight\il{Berber!Tamazight}/Amazigh\il{Berber!Amazigh|see{Berber!Tamazight}}'', and occasionally as ``Berbisch\il{Berbisch|see{Berber}}''. The designation ``\ili{Berber}'' is regarded as pejorative by many members of the community, who therefore oppose its use. The variety spoken by Hilal and Lamin is Tarifit\il{Berber!Tarifit}, i.e., Riffian\il{Riffian|see{Berber!Tarifit}}. Like in other varieties of the language, Tarifit\il{Berber!Tarifit} speakers refer to their language simply as ``Tamazight\il{Berber!Tamazight}'' in their own language. In that sense, ``Tamazight\il{Berber!Tamazight}'' (or [the masculine form] ``Amazigh\il{Berber!Amazigh|see{Berber!Tamazight}}'', when used as an adjective in \ili{English}) is equivalent to what is often referred to as ``\ili{Berber}'' in academic literature as well as in popular discourse. Note that the label ``Tamazight\il{Berber!Tamazight}'' is ambiguous in that it refers to the larger family of linguistic varieties (``\ili{Berber}''), but also in a more specific sense to a smaller group of varieties used mainly in parts of the Central Atlas mountains. Therefore, we decided for the use of Tarifit\il{Berber!Tarifit}, but use other terms according to speakers usage in the conversation, when we quote or paraphrase them.} but has an excellent command of Moroccan \ili{Darija}, too (Hilal [00:06:28--00:06:40; 00:16:04--00:16:10]). 

The participating researchers' role was to keep the conversation going. The other three participants were asked to tell us about their experiences. Of them, only Marcus -- a linguistics student -- had received some information prior to the meeting concerning the fact that we were interested in topics related to experiences of discrimination, but also in the way these would be talked about. He had been instructed to steer the topic towards experiences of discrimination, either by contributing examples himself or by asking the Tarifit\il{Berber!Tarifit} speakers about theirs. The three non-researcher participants were explicitly invited to discuss the key theme among themselves and to ask each other questions whenever they did not understand or simply wanted to comment.

\ili{German} was the language used in that conversation, which lasted about two hours in total. Our objective was to create data that could be analyzed on two different levels. On the one hand, we were interested in hearing these young men's stories. Are there recurrent topics? What matters to them? How do they deal with experiences of exclusion and othering that -- we assume -- all of them will have made? On the other hand, we wanted to obtain data that would tell us about power in language by looking at the conversation situation from a linguistic point of view. For that reason, it was important to us to have a mixed group, with two researchers of different language backgrounds, as well as ``native'' speakers of \ili{German} and Tarifit\il{Berber!Tarifit} among the participants. How would these differences play out in terms of accommodation and adjustment of one's behavior in a group with different competencies in \ili{German}? Do their ways of speaking attest to underlying expectations concerning language skills and communicative competence of the others? 

Key sequences from the conversation include an episode of a police check during a holiday stay in Nador, Hilal's northern Moroccan hometown. Hilal experienced what he considers inappropriate communicative behavior by a policeman whom he perceived as a linguistic and ethnic ingroup member. He describes the incident, which is then commented on by Lamin. Moreover, the use of derogatory ethnic or social group labels is worth looking at. This ranges from Marcus's explanations concerning examples of pejorative terms that were used -- insultingly or provocatively -- in his presence and challenged him as a \ili{German} individual to the subverted use of such labels among members of an ingroup. 

Why do we focus on sequences in which the interlocutors deal with (verbal) experiences that they perceived as discriminatory? The reason is that, in these cases, different perceptions of such situations are likely to trigger strong emotional reactions and possibly dissent. In such situations, communication is no longer perceived as adequately aligned according to convergence principles of CAT. When speakers gain the impression that adjustment mechanisms are mal-functioning or, more significantly, wilfully violated, they will develop a strong sense of difference, possibly even adversity. The discussion of such experiences in a group conversation is an important stimulus for creating a metapragmatic conversation. Examples include instances in which Hilal and Lamin evaluated particular situations differently. 

\section{Conceptual tools for the analysis of nonaccommodation}\label{CH1.1.ss.3}

\subsection{Interpretive strategies between CDA and CAT}\label{CH1.1.ss.3.1}

Our attempt at analyzing this conversation is informed by two rather different approaches. On the one hand, we rely on input from CAT; on the other hand, we try to understand what happened during the conversation through an interpretation that draws on the tenets of critical discourse analysis\is{critical discourse analysis (CDA)} (CDA) as proposed originally by Norman \citet{fairclough1995,fairclough2013} and Ruth \citet{wodak_meyer2009}. Without aiming at a full-fledged contrastive evaluation of these perspectives, we do hope to contribute to a bigger question that challenges both, albeit in different ways. In terms of its central concepts, CAT draws on social psychology: human categorisation and group membership are crucial; so are notions of asymmetry, bias and linguistic dominance. All of these are necessary to understand the dynamics of the conversation we held (both in terms of contents and manner). It is important to bear in mind that CAT as a theoretical framework is not a uniform, entirely formalized theory. It may, at times, be challenging to identify its tenets without making it look overly categorical. As a matter of fact, as a theoretical framework, CAT is fairly adaptive and takes into account the significance of sociohistorical contexts, social settings, cultural contexts with their possibly malleable group affiliations and construal mechanisms of identity that allow for multiple group membership. Nonetheless, it is also important to delineate CAT as a conceptual approach from an entirely open-ended deconstruction or critical approach in sociolinguistics. 

\begin{sloppypar}
A key dimension is the question of to what extent individual identity is perceived (both by the individuals ``under investigation'' as well as by those analyzing the data) as linked to membership in a group that exists outside the categories established in and through discourse. CAT suggests that social communication is dynamic and that interlocutors’ identities and group memberships may shift and vary in salience both across interactions and within a conversation, (e.g., depending on the topic of the conversation, interpersonal commu- nication can shift to the intergroup level). Central to CAT are notions of ingroup- versus outgroup-communication implying that groups exist as fairly stable, easily identifiable and clear cut bounded categories. How does this relate to notions of fluidity of such categories, fuzzy boundaries, changing loyalties? 
\end{sloppypar}

CAT and CDA are comprehensive theoretical frameworks. Both pay attention to a broad range of possible features and contextual aspects affecting how meaningful communication evolves in practical interaction. CAT tends to place more emphasis on how communication works as a collaborative effort (or where it fails to do so) in actual dialogues and conversations. Critical discourse analysis\is{critical discourse analysis (CDA)} highlights power asymmetries established, but also possibly subverted, in actual conversation. 

There is a broad overlap between both frameworks, enabling us as a team of authors to communicate despite different backgrounds in terms of academic disciplines and, accordingly, distinct familiarity with the frameworks and their methodological implications. Our joint venture has, therefore, also been an exercise in testing the compatibilities of both approaches or -- where we would find significant differences -- the potential benefit of drawing on both contrastively. 

Both approaches are deeply data-driven and empirical. All three authors did close readings and re-readings of the conversation, before discussing the material and our respective interpretations. In these discussions, we did not always agree on how to interpret what was going on during the recorded conversation.

The ``double-faced'' character of language as, on the one hand, our most significant tool and, on the other hand, something that imposes itself upon us and guides us towards different ways of thinking is mirrored by the question of what a relevant group is (as a conceptual, analytical category in CAT, but also beyond). Groups exist as cognitive models, and as such, they are strong, lasting, and seemingly clear-cut (despite being commonly contested). Human beings behave according to their ideas of what possible categories exist and matter. In that sense, accommodation and nonaccommodation are clear mechanisms targeting discernible entities.

With the appearance of what \citet{eckert2012} labels as the third wave in sociolinguistics, attention was drawn increasingly to mechanisms of stylization and the indexicality of linguistic markers, accompanied by the recognition of how significant language ideologies are in everyday verbal practices \citep{kroskrity2010,spitzmuller_busch_flubacher2021}. \citet[6--7]{levon_mendes2016} warn against an instrumental notion of agency (= ``doing linguistic identity'') while recognizing the significance of stance, agency, and the opportunities they afford to speakers to position themselves through ways of speaking. This implies that our linguistic behavior constitutes group membership rather than being a direct function of it (for an in-depth discussion of language and agency, see also \citealt{duranti2005}). Where these approaches and CAT meet is the acknowledgement of the significance of human agency, always bearing in mind that agency does not necessarily mean premeditated behavior. CAT concerns itself specifically with how speakers adjust to one another, be it consciously or not. While CAT is not a formal theory, but an adaptive theoretical framework, it offers a somewhat standard conceptual and terminological toolset. These notions have gained currency among those working within the framework, but others may not be equally familiar, which is why we deem it important to dedicate some space to preparing the theoretical ground in terms of CAT before delving into the concrete examples drawn from the recorded group conversation. 

\subsection{Speech adjustments and failure to accommodate}\label{CH1.1.ss.3.2}

Sociolinguistic work has established that the way people interact with each other and the way they adjust their speech depends on ``who speaks what language to whom and when'' \citep{fishman1965}. Social variables interacting with each other in complex ways play a significant role, as they shape speech, language and communication (e.g., \citealt{chambers2003,eckert2018,ferguson1959,ferguson1975,fishman1965,goffman1964}). As \citet[133]{goffman1964} stated, ``[i]t hardly seems possible to name a social variable that doesn't show up and have its little systematic effect upon speech behavior: age, sex, class, caste, country of origin, generation, region, schooling; cultural cognitive assumptions; \isi{bilingualism} and so forth''. Group membership is one of those social factors that does not only involve determinants, such as age, \isi{gender} or ethnicity, but also, as \citet[68]{fishman1965} argued, is closely related to socio-psychological variables. Depending on the situation, people may belong to different groups changing their register or switching to another language accordingly. The desire to belong to a group or gain acceptance from it also influences speech behavior. Additionally, factors such as interactants' expectations, motivations and perceptions affect their speech and communicative adjustments \citep{beebe_giles1984,giles_coupland_coupland1991}. This implies that interactants not only may adjust their communicative behavior to their interlocutors unconsciously and automatically but also consciously and intentionally (\citealt{dragojevic_gasiorek_giles2016}; see also \citealt{gasiorek2016b} for a concise overview of research investigating individuals' awareness of the adjustments they make). 

One of the assertions of CAT is that speakers' adjustments are guided by cognitive and affective motives \citep{gallois_ogay_giles2005}. Cognitive motives concern comprehension and may be associated with needs and desires to communicate effectively: Speakers may converge to the recipients' speech characteristics or communicative behavior to facilitate comprehension \citep{dragojevic_gasiorek_giles2016,gallois_ogay_giles2005}. For instance, in intergroup communication between native and non-native speakers, native speakers may simplify their vocabulary and syntax \citep{petrou_dragojevic2024,zuengler1991} to aid understanding. Similarly, speakers may maintain or diverge from their interlocutors to enhance communication \citep{dragojevic_gasiorek_giles2016,gallois_ogay_giles2005}. For example, a speaker may diverge from an interlocutor who speaks overtly fast by speaking slowly to urge them to do the same (\citealt{brown_giles_thakerar1985} as cited in \citealt{dragojevic_gasiorek_giles2016}). Divergence may also serve as a strategy to signal the need for the interlocutor to adjust their speech to better suit the context \citep{gallois_ogay_giles2005}. Additionally, it may be intentionally applied to make a conversation problematic or to deprive (some of the) interlocutors of understanding, for example, by engaging in code-switching \citep{gasiorek2016a,dragojevic_gasiorek_giles2016} or by changing the code \citep{petrou2023}.
As mentioned earlier, CAT posits that, in addition to cognitive motives, adjustments are also made based on people's affective motives. These are related to negotiating personal and social identities. Particularly, CAT proposes that people regulate social distances through communication in order to establish positive personal and social identities and reach social approval \citep{dragojevic_gasiorek_giles2016,giles_coupland_coupland1991}. This can be accomplished by converging to one's communicative behavior ``to appear more similar and thus more likable'' \citep[126]{gallois_ogay_giles2005}. People may also seek to differentiate themselves from others positively. This may happen when they compare themselves to others and find, or even create, dimensions that they perceive as better than those of others \citep{tajfel1974}. Distinctiveness is accomplished by maintaining or diverging from their conversational partner (\citealt[126]{gallois_ogay_giles2005}, \citealt{dragojevic_gasiorek_giles2016}). For instance, native speakers may switch their register from standard to a dialectal variety when communicating with non-native speakers, not only to make communication problematic but also to mark their distinctiveness from outgroup members (see example 7 in \citealt[140]{petrou2023}). Another example of this ``psycholinguistic distinctiveness'' \citep{giles_bourhis_taylor1977} provides the sociolinguistic situation in northern Cyprus. Cypriot Turks maintained their Cypriot \ili{Turkish} \isi{dialect}, which is now a prestigious variety, to differentiate themselves from settlers from Turkey in the northern part of Cyprus in light of cultural differences \citep{kizilyurek_gautier-kizilyurek2004,petrou2023}.

CAT approaches communicative adjustments not only from the perspective of the speaker but also from the perspective of the recipient. Recipients experience communicative adjustments as \textit{accommodative} when they perceive them as appropriately adjusted \textit{relative} to their needs and as \textit{nonaccommodative} when they experience them as inappropriately adjusted \textit{relative} to their needs (for a recent study, see \citealt{petrou_dragojevic2024}). Recipients may perceive nonaccommodative interactions and speakers as \textit{overaccommodative}, i.e., they are perceived to exceed the required level of adjustment, as \textit{underaccommodative}, i.e., they are perceived to undershoot the required level of adjustment, or as \textit{counteraccommodative}, i.e., they are experienced as to be ``psychologically dissociative'' \citep[32]{coupland_coupland_giles_henwood1988}. 
While accommodative adjustments are positively evaluated, nonaccommodative adjustments are negatively evaluated. As discussed in Section \ref{ch.1.2.ss.4}, negative evaluations are associated with ascribed motives \citep{gasiorek_giles2012}. In intergroup communication, recipients are more likely to ascribe negative motives to conversational partners who are seen as outgroup members and positive motives to conversational partners who are ingroup members \citep{gasiorek2016a}. Most interactions with the ``sense of ``us'' and ``them'''' \citep[308]{gallois_watson_giles2018} are intergroup. In fact, interpersonal communication may quickly switch into intergroup (e.g., \citealt{gallois_watson_giles2018}) or ``be infused with group identities'' \citep[30]{dragojevic_giles2014}. In our episodes, the ingroup component is strong, and the sense of ``we'' and ``they'' prevails.
With these conceptual tools in mind, we can delve into the first case, an episode that relates Hilal's experience with a person that he considered a member of the same group as himself, on the basis of their assumed common linguistic background as speakers of Tarifit\il{Berber!Tarifit}. 

\section{Speaking about nonaccommodation: a conversation}\label{CH1.1.ss.4}

\subsection{Angered by nonaccommodation: a police encounter in the Rif}\label{CH1.1.ss.4.1}

What follows in this section is an example of divergence or nonaccommodation within one's own group. There are at least two aspects that we deem interesting about this particular incidence. 

First, it serves to scrutinize notions of boundaries between groups. This question is of immediate significance for our chapter, because one of its objectives is to contrast views that need to assume linguistic communities as fairly stable categories and, thereby relatively strict notions of group membership with other views that assume more fluid boundaries, agency in sociolinguistic choices and actual construal of linguistic communities by means of language (use), rather than their extra-linguistic existence pre-determined through factors outside language.

Second, the episode is of particular interest because two of the participants in the group discussion that we arranged, Hilal and Lamin, differ as to how they assess the situation. Hilal, who narrates this episode as a personal experience during a holiday stay with the family in the Moroccan Rif region, felt discriminated against by a member of what he perceives as his own speech community, North African Tarifit\il{Berber!Tarifit} speakers. Lamin's attitude is more distanced; he is not personally involved and is calm, in contrast to Hilal, who was noticeably angry when talking about the incident. The episode is, therefore, not only interesting for what it reports on but also for the situation it brings about in the group interview. 

Hilal shares an incident that happened to him in his northern Moroccan hometown Nador. Nador's population is mostly Tarifit\il{Berber!Tarifit}-speaking -- like Hilal -- and has grown rapidly over the past fifty years, largely due to migration from rural areas in \isi{Morocco}'s northeastern L'Oriental region. Hilal told us that a policeman from Nador stopped his car. The conversation between both took place in Moroccan Arabic\il{Arabic!Moroccan}, known as \ili{Darija}, as the policeman did not accommodate Hilal, who initially communicated in his ``mother tongue'', a variety of \ili{Berber} called Tarifit\il{Berber!Tarifit}. The policeman diverged from him by communicating solely in \ili{Darija}. As we do not know who initiated the dyadic conversation, it is possible that it was Hilal who initially diverged: probably, the policeman addressed him in \ili{Darija}, and Hilal replied in Tarifit\il{Berber!Tarifit}. Regardless of which of the two scenarios transpired, the individual who adjusted his communicative behavior towards their conversational partner was Hilal: 
\begin{quote}
\textit{Und dann habe ich ihm gesagt ja, auf Arabisch die ganze Zeit. Am Anfang habe ich meine Muttersprache \textnormal{[}gesprochen\textnormal{]}, weil ich bin so, weil ich rede meine Muttersprache, ich habe Tarifit\il{Berber!Tarifit} gesprochen, er antwortete \textnormal{[}dar\textnormal{]}auf mit Arabisch}. 

`And then I said to him, yes, in \ili{Arabic} the whole time. At first it was my mother tongue that I [used], because I am like that, because I talk in my mother tongue, I spoke Tarifit\il{Berber!Tarifit}; he answered in \ili{Arabic}.' (Hilal [00:15:55--00:16:18])\footnote{All translations by the authors.}
\end{quote}

As the extract above shows, Hilal said that he initially spoke in Tarifit\il{Berber!Tarifit} ``because I speak my mother tongue''. We could posit that his \textit{initial orientation} (\citealt{dragojevic_gasiorek_giles2016}), directed by his perceived social identities regarding the policeman, was that the latter would be a speaker of Tarifit\il{Berber!Tarifit} like Hilal himself. When he was asked by Axel how he knew that the policeman was able to speak Tarifit\il{Berber!Tarifit}, Hilal replied that he knew for sure that the policeman spoke Tarifit\il{Berber!Tarifit} when he communicated in Tarifit\il{Berber!Tarifit} with another policeman who arrived at the scene later: 
\begin{quote}
\textit{Kam sein Kumpel auf einmal kommunizieren sie auf Tarifit\il{Berber!Tarifit}. Und das war für mich Diskriminierung. Habe ich gesagt: „Warum hast du nicht mit mir jetzt auf meine Muttersprache gesprochen?{“} Er sagte zu mir: „Nein es gibt Gesetze und wir müssten uns an Regeln, nicht Gesetze, Regeln.“}

`Then came his mate, and suddenly they communicate in Tarifit\il{Berber!Tarifit}. And that was clearly discrimination in my view. So, I said, ``Why didn't you talk to me in my mother tongue just now?'' He said to me, ``No, there are laws and we need to abide by the rules, not laws, rules.{''}' (Hilal [00:16:35--00:17:14])
\end{quote}

Perhaps there were indeed instructions or rules that policemen should use the \ili{Arabic} language during their duty interactions with civilians. It is also possible that the policeman used \ili{Arabic} to indicate that he belonged to or had a preference for the \ili{Arabic}-speaking group in order to (positively) differentiate himself from Hilal in terms of status -- him being a law enforcement agent, Hilal, a mere civilian. We posit that the conscious use of \ili{Darija} served different affective functions: it regulated the social distance between the two -- \ili{Arabic} reinforced their different status, the effect of which could be the discouragement of Hilal from communicating in a way that might have been inappropriate in the given context. 

Although Hilal had a good command of Moroccan Arabic\il{Arabic!Moroccan}, he experienced the use of \ili{Darija} as ``discrimination'', a clear case of counteraccommodation directed against him. On the other hand, Lamin, who grew up in a mixed family, one parent being \ili{Arabic}, the other Tarifit\il{Berber!Tarifit}-speaking, perceived the use of \ili{Darija} differently. Just like Hilal, he was convinced that the policeman could be expected to be a speaker of Tarifit\il{Berber!Tarifit} in the local setting in the Rif region. However, Lamin gives another reason for why the policeman spoke \ili{Darija} instead of Tarifit\il{Berber!Tarifit}. He considered him simply more open-minded about the choice of language: \textit{Sie sind offener halt} `They [referring to the policemen] are simply more open' (Hilal [00:17:43--00:17:46]). Then, Hamza skeptically noted that if the policeman had been open-minded, he would have talked to Hilal in Tarifit\il{Berber!Tarifit}. Hilal agreed fervently and added that the policemen intended to provoke him. He assumed that this was due to instructions from their superiors. This indicates that Hilal inferred negative motives for the language choice of the policeman. To Hilal's inferred motives, Lamin responded with astonishment: \textit{Warte, er ist \ili{Berber}, er redet mit seinem Kollegen Berbisch\il{Berbisch|see{Berber}}, und dann provoziert er einen \ili{Berber}?} `Wait, he is \ili{Berber} himself, he talks \ili{Berber} with his colleague, and at the same time provokes a \ili{Berber}?' (Lamin [00:18:35--00:18:43]).

What can we learn from this episode? Overall, the situation does not seem particularly noteworthy at first sight: a policeman stopping a car whose driver did not follow a valid traffic rule. There is more to it, though, because of how the dialogue unfolds. We need to bear in mind that all we know about the incident, we have learned through Hilal's narration. Yet, even with this likely bias in mind, there are several points we can draw. Some of them concern the construal and perception of speech communities and group identities within the episode narrated by Hilal; others concern the meta-pragmatic assessments and the contrast between Hilal and Lamin in this regard. Let us look at these in turn.

It is important to ask who initially set the scene. Hilal was stopped by the policeman, who must have addressed him first -- presumably in \ili{Darija} -- so we can assume that it was Hilal who ``switches'' first. While Hilal blames the policeman for not speaking Tarifit\il{Berber!Tarifit} to him, we need to bear in mind that Hilal first pushed for a ``language reset'', changing to Tarifit\il{Berber!Tarifit} when the policeman addressed him. If Hilal had reacted in \ili{Darija}, the course of the conversation and the outcome -- a fine (?) -- would have been largely the same, albeit with no indication that the linguistic construal of the respective identities, language choices or other verbal behavior were an issue. Had Hilal reacted less defiantly in \ili{Darija}, he would still have felt discriminated against. In fact, that is exactly how things turned out: eventually, he gave in and used \ili{Darija}, complaining in hindsight that everyone has to talk in \ili{Darija} and not Tarifit\il{Berber!Tarifit}. 

There is an interesting twist to this. Hilal, who would appear as more prone to navigate different registers, alluding to group membership in different ways, turns out to rely on a much stricter categorization scheme than the policeman. For Hilal, being a speaker of Tarifit\il{Berber!Tarifit} is a central feature of people from his native region. In practical terms, Tarifit\il{Berber!Tarifit} is used widely, and Hilal can certainly expect a local policeman to be conversant in that language -- even before this is confirmed by the policeman talking to a colleague in Tarifit\il{Berber!Tarifit}. Identification with a speech community builds an ingroup/outgroup scenario that is taken for granted, perceived as natural and not drawn into question by Hilal. Irrespective of the policeman's rigorous attitude, his norm-enforcing function and his entire presence linked to state authority, he can be argued to be much more fluid in terms of linguistic \isi{belonging} than it appears at first. 

Of course, the policeman may have chosen \ili{Arabic} in order to render the encounter formal, distancing himself from Hilal through this nonaccommodative kind of behavior. This is Hilal's assumption, which he clings adamantly to. It may be the case that Hilal misjudged him in assuming that everyone would rely on their native language as the central identity-construing property, and, consequently, the primary and dominant language to choose in any encounter with an unknown interlocutor in public space. After all, it is perfectly conceivable that the policeman simply used one of several codes, here \ili{Darija}, automatically for no particular reason, when he addressed a driver who had infringed a traffic rule. Ultimately, considerable uncertainty remains in this respect.

Communication-related research has suggested that people's communicative behaviors can carry varying interpretations. This is because different recipients may perceive the same communicative behavior differently (e.g., \citealt{dragojevic_gasiorek_giles2016}; see also \citealt{gallois_weatherall_giles2016}). People's perceptions are partly influenced by factors including sociocultural norms and cultural communicative repertoires. As the literature on CAT points out, nonaccommodative strategies are typically defined by recipients' perceptions \citep{gasiorek2016a}. How recipients react to nonaccommodative strategies is determined by the way they interpret the speakers' motives \citep{gasiorek_giles2012}.

In this regard, it is noteworthy that Hilal and Lamin assess the situation differently. It is impossible to tell who is right or wrong; that is simply not a substantial point here. More interesting, as it turns out, is what both read into the situation and how they discussed it during our conversation. According to Lamin, the policeman did not necessarily have any intention to discriminate against Hilal through his choice of language. Hilal's and Lamin's differing views may appear to align with their respective stance towards matters of \ili{Berber} identity, with Hilal perceiving the situation as threatening to his ethnic identity (\citealt{bourhis_giles1977} for a discussion of relevant mechanisms), and Lamin not sharing this perception, perhaps due to his mixed background. 

On a meta-pragmatic level, we may wonder what Lamin intended to achieve when he made those few remarks that he contributed after Hilal had told his story. Are they intended to soothe Hilal while depicting himself as a mediator? Or are they a power mechanism, positioning himself as the more worldly-wise, rational person of the two during our group conversation? Given that Lamin grew up in a family of mixed language backgrounds, this may have been an attempt at portraying himself as someone who defies straightforward ethnic categorization and handles situations of potential conflict better. While Lamin and Hilal differ in how they interpret the police episode, both tell us many stories about incidents in Germany that they experience in a very similar way. 

\subsection{Microagression and verbal assault: conventionalized nonaccommodation}\label{CH1.1.ss.4.2}

In the Moroccan scenario, the relevance of \isi{belonging} to a particular group and the need to establish which construal of outgroup versus ingroup communication participants in the discourse rely on Hilal's encounter with the policemen (and the differing interpretations by Hilal and Lamin). In contrast, the situation that Hilal and Lamin experience as male North African migrants in Germany appears more obvious. In a sense, one could logically suspect that any kind of boundary work establishing significant categories is somewhat less significant under the circumstances. When social categories such as ``member of the host community'' versus ``(im)migrant'' are taken for granted, it might look like there was not much need to establish them discursively. However, even when such categories appear objectively, fairly clearly, and significantly, individuals keep drawing attention to them in discourse. Their obviousness does not render discriminatory acts towards marginalized groups less common. To the contrary, discrimination based on a categorization of people into -- arguably obvious -- groups and communities with alleged social identities is common. 

Hilal experienced discrimination in \isi{Morocco}, as he told us, but certainly in Germany, as well -- not only from members of the host community but also from individuals who -- just like him -- had immigrated to Germany. In particular, he recounted an incident (Hilal [00:02:50--00:03:10; 00:07:17--00:10:00]) in which a young man from a central European country asked him to go ``back to Africa'' (Hilal [00:03:09]). Such experiences cause emotional strain: Hilal described that he felt attacked (\textit{angegriffen}; Hilal [1:41:13--1:42:10]) and voiced that ``after a month this was still in my head'' (Hilal [01:41:34--01:41:44]). Such kind of discriminatory behavior -- a verbal assault -- is definitely aggressive. It is hard to imagine how such an offense should occur unintentionally. 

Hilal's reaction is interesting in that it contains more than just anger. Hilal explained in some detail (Hilal, with short interjections by Lamin and Hamza [01:40:40--01:45:45]) that it was his dream to be in Germany and when someone discriminates against him on the grounds of ``not \isi{belonging}'' (Hilal [01:42:30--01:42:52), he feels greatly disappointed and dejected. He points out to us that how he feels about or responds to discrimination depends on various factors. These include the way people verbalize their discriminatory attitude. It also matters who discriminates -- are these immigrants themselves, or not? Moreover, Hilal mentions mistakes on his part as a reason that could trigger discriminatory offenses against him, insinuating that there could be reasons that could at least partly justify discriminatory behavior (Hilal [01:43:50--01:44:40]). Where such ``reasons'' are absent, Hilal states, discriminatory behavior makes him sick: 
\begin{quote}
\textit{Wenn jemand jetzt zum Beispiel mich nur wegen AUssehen und der Sprache oder Hautfarbe diskriminiert, da, da bin ich krank.}

`If someone discriminates, for example, based only on the appearance or the language or the skin color, then I am, then I am sick.' (Hilal [01:45:01--01:45:14]) 
\end{quote}
Lamin agreed and added: 
\begin{quote}
\textit{Manchmal musst du nicht mal was machen, also als Ausländer mit so'ner Hautfarbe, auch wenn du … also meist sind das die Älteren, wenn du auch nur so vor Supermarkt, zum Beispiel, guckt, sieht dich, das Erste, was sie macht, die Tasche nach vorn.}

`Sometimes you do not have to do anything; that is, as a foreigner with such a skin color… most of the times it is the elderly, when you are in the supermarket for example, she looks, sees you, the first thing she does, the bag to the front.' (Lamin [01:45:14--01:45:32])
\end{quote}

To this, Hilal reacted vividly, stating that he experienced this every day. To this incidence, we want to linger for several reasons. First, because people do not only communicate verbally; communication is also nonverbal. Nonverbal cues, such as clothing and appearance, inform interactants about their interlocutors' group memberships and social identities (e.g., \citealt{giles_maass2016,keblusek_giles_maass2017}). What Hilal and Lamin regularly experience due to their physical appearance is in line with findings of other research dealing with experiences reported by immigrants/non-native speakers living in Germany (\citealt{petrou2023,petrou_dragojevic2024}).

Both Hilal and Lamin -- as members of the same group of male North African migrants -- perceive and evaluate the action of moving the handbag in front in the same way, very negatively. Interestingly, members of the host community doing this action did not communicate anything to Hilal and Lamin straightforwardly (for perceptions that were not directly communicated but inferred, see \citealt{petrou_dragojevic2024}). Yet, Hilal and Lamin coincide in how they feel about such situations, as being perceived as a potential threat by outgroup members. Through the lenses of CAT, this could be another incident experienced as nonaccommodative, specifically nonaccommodative in terms of emotional expression or even as counteraccommodative. It can be concluded that due to prejudice based on appearance, members of the host community underaccommodate the affective needs of people with darker skin color influenced by stereotypes they hold about groups of people having specific characteristics. This kind of communication -- even if in this case, it is indirect and nonverbal -- that ignores or hurts the feelings of the ``other'' can be perceived as psychologically dissociative and hostile, thus as counteraccommodative. 

Other explanations for what Hilal observed -- an elderly woman quickly grasping her handbag -- are possible. She may have acted for some reason other than fear of being pick-pocketed by a foreign-looking young man. It might simply have been coincidental. Nevertheless, the significance of such incidents should not be downplayed. Both Hilal and Lamin give several examples and stress that it happens every day. The apparent mismatch between judging such incidents as banal, on the one hand, and their frequent experience and perception of hostility, on the other hand, are reminiscent of the debates around microaggressions. 

Research on microaggressions has been criticized on the grounds of its subject matter being too evasive. There is a genuine dilemma at work in that case. When those who are affected by microaggressions on the basis of \isi{belonging} to a particular group, e.g., migrants, point this out to non-migrants, they are typically met with skepticism. This is an inherent feature of microaggressions. Outsiders, i.e., members of the larger host community, tend to assume that such perceived aggressions in the form of verbal assaults or invalidations are being misunderstood, exaggerated and blown out of proportion by those affected. Often, this is accompanied by the assumption that the latter are overly sensitive. 

The dilemma is that those relying on stereotypical notions, even if inadvertently, are not usually questioned about their actions or utterances. Lamin and Hilal both tell us about such experiences. Both are aware of many situations in which they felt they were being perceived as a potential threat by outgroup members. These situations are usually not drawn to the level of conscious consideration by the outsiders, in other words, the perpetrators, let alone being addressed explicitly. If, exceptionally, they are, the exculpatory mechanisms --- as mentioned above --- serve to save one's face. It is much more common for matters to remain implicit and unaddressed. For those targeted by such microaggressions (including entirely unintended acts or utterances that demonstrate implicit categorizations and notions of alterity) this means that they have to deal with the situation choosing from among a small range of options. Possible reactions include downplaying it, ignoring it, dealing with one's anger silently -- but only very marginally confronting the other about it. What's more, even if such situations were to be addressed explicitly between those involved, with an intent to clarify matters, this would not undo the boundary between these groups that are assumed as real, natural distinctions by their respective members. If anything, it would enhance them. 

In summary, we see that boundary work, often rather inconspicuous yet powerful in establishing societal fault lines, is also at work in settings where the differentiation between groups appears as obvious as in the case of North African migrants in Germany. What constitutes in- and outgroup communication is not a given; it follows pervasive mechanisms in communication and interaction at large. Hilal and Lamin describe this quandary as deeply saddening -- sickening, even. The subtle incidents like someone grasping their bag in the supermarket leave little room for addressing this; being coerced to remain silent about it explains what Hilal and Lamin describe as saddening or sickening. With explicit aggression, the kind of verbal assault that must necessarily be understood as intentionally offensive, matters look a bit different. Looking at named categories of groups, including those labeled in very derogatory terms, is another component of what constitutes experiences of discrimination. We turn to them in the following section.

\subsection{The substance of verbal offense: ethnophaulisms and endearment}\label{CH1.1.ss.4.3}

\begin{sloppypar}
As indicated above, inter-group communication triggers social comparison driven by the desire to differentiate oneself positively from outgroups. Ingroup favoritism may lead to ascribing unfavorable traits and characterizations to the outgroup. In fact, \citet{keblusek_giles_maass2017} speak about commonly occurring derogatory labels or ethnophaulisms, which they describe as linguistic markers of hostile \isi{attitudes} toward the outgroup. Derogatory labels are used to belittle and attack the reputation of outgroups. This is commonly achieved by drawing on features associated with the outgroup that are perceived as negative, in contrast to highlighting aspects deemed positive about one's own social identity, for example, in terms of religious beliefs. For instance, Christians and \isi{Muslims} may apply terms such as ``heathens and infidels'' to one another, depicting the respective other group of people as nonbelievers \citep[637]{keblusek_giles_maass2017}. Conceptual metaphor and metonymy are common mechanisms through which words acquire new meanings \citep{murphy2010}. These, also serve as language tools typically used to portray ingroups' perceptions and evaluations of outgroups. Among them food and animal expressions are copious \citep{keblusek_giles_maass2017,haslam_holland_stratemeyer2016}.\footnote{%
Interestingly, food and animals are also used in argot or so-called ``secret languages'' where ingroups engage in a deliberate change of meaning of words so that they won't be understood by outgroups by applying tools such as metaphors and metonymies -- often in correlation with components such as comparison and experience \citep{petrou2020}. For instance, in \ili{Turkish} argot \textit{angoraki} (< gr. `small cucumber; cucumber') denotes `stupid; stupidly naïve', \textit{kuzu} (tr. `lamb') denotes for example an `inexperienced or immature person, young girl' \citep{petrou2020}.} For example, Americans call \ili{French} people ``frogs'', \ili{French} call \ili{English} people \textit{les rosbifs} `the roast beefs' and Italians ``\textit{les macaronis}'' (Data from an interview with a \ili{French}-speaking sojourner participant in Germany; \citeauthor{petrou2023}, \citeyear{petrou2023}). 
\end{sloppypar}
 
It did not take much effort to prompt the participants in our group interview to tell us some such labels marking intergroup relations, used for members of the outgroup. Marcus explains his reaction to being called \textit{Kartoffel} `potato' or \textit{Alman} (original \ili{Arabic} `\ili{German}', now borrowed into migrants' varieties of \ili{German}), or overhearing people using these terms that evoke negatively valued attributes associated with stereotypical Germans, such as being rigid, narrow-minded, smart-alecky, humorless.

While ethnophaulisms in the narrow sense are coined by one group to refer to another and portray its members in a negative light, a different type of abusive language criticizes what is thought of as a disloyal attitude towards the ingroup. In France, Africans who show a decidedly positive cultural or political stance towards France are sometimes labeled as ``bounty'' by other Africans, based on a coconut-filled, chocolate-glazed candy bar with that name: black on the outside, but white inside (\ili{French}-speaking sojourner, interview data, \citealt{petrou2023}). \citet[639]{keblusek_giles_maass2017} mention similar cases from North America: ``coconut'' for Black individuals, ``apple'' for Native Americans when individuals would associate themselves culturally with the American mainstream, while other members of their alleged ingroup use such labels to express their disapproval of what they regard as cultural betrayal. 

We have not really noted any such terms during the interview, and in fact, we have not mentioned them in informal conversations before or after the recording device was switched on. What we did note, however, were several labels that could carry rather different, seemingly contradictory meanings. Connotations of the designation \textit{Araber} (\ili{German} for `Arab') appear to depend largely on the context in which it is used. There is nothing inherent in the word that would make it offensive as such. In many instances, Hilal, who is usually adamant about being an Amazigh\il{Berber!Amazigh|see{Berber!Tamazight}} in terms of culture, ethnicity and language, would not object to being called an Arab. In other contexts, Hilal experienced being called an Arab as inappropriate. Pejorative meaning is not an inherent property of the word. It is the assumption Hilal makes about why the label is used to refer to him that affects his reaction to it. Whether it matters or not depends on the intention that Hilal ascribes to those using it. Accordingly, it may work as a hurtful or provocative label, or end up being a moot offense. 

\begin{sloppypar}
Lamin also talks about an ambivalent term. In Tarifit\il{Berber!Tarifit} and other Amazigh\il{Berber!Amazigh|see{Berber!Tamazight}} varieties \textit{aberkan} means `black'. Lamin points to his dark complexion and explains how he was called \textit{aberkan} as a young child in school (Lamin [00:23:55--00:24:40]). The label remained and turned into a nickname for him. As such, the term was not used in a demeaning manner despite the negative connotation it carries. This affected him as a child, which he remembers quite vividly. When we talked about how certain terms referring to ethnic groups, social identities or personal properties of a person can turn into abusive language and attain a function as swear words, he would refer to skin color in particular. 
\end{sloppypar}

Stigmatized groups may use social creativity in various ways to positively influence the way they are viewed. This can be achieved by causing revaluation of negative labels \citep{galinsky_hall_cuddy2013,galinsky_wang_whitson_anicich_hugenberg_bodenhausen2013}. A well-cited example of this process is the ``Black is Beautiful'' movement, which was an attempt at changing the negative connotation of Blackness \citep{galinsky_hall_cuddy2013,keblusek_giles_maass2017}. Another act of reappropriation is the use of derogative labels by the stigmatized or \isi{minority} groups themselves, e.g., the use of the term ``queer'' by some, but by no means all, LGBT+ people. 

The playful use of \textit{Kartoffel} or \textit{Alman} (briefly discussed by the participants during the conversation, [00:57:54--00:58:06; 00:01:03--00:01:58; 00:10:00--00:11:18]) in comedy shows, initially and mainly by verbal artists with ties to North Africa or the Middle East, appears to catalyze such use even as an ingroup term at least in a satirical manner or self-mockingly among Germans. (On a side note, Lamin jokingly mentioned Rap and Comedy as the two arenas where being North African might be grounds for positive discrimination in Germany [00:43:23].)

Whether an insult works or fails is context-dependent. This implies, of course, that these terms can depend on which two (or more) people use them. Hilal and Lamin greet each other with a resounding \textit{Hi, du Flüchtling} `Hi, you ol' refugee' [00:57:23--00:57:53]), which works as a term of endearment among these North African men. The \ili{German} word \textit{Flüchtling} has become connotated negatively over a relatively short time-span and is often avoided now, using \textit{Geflüchtete*r}, a nominalized participle, instead. The non-offensive use of originally offensive terms can be instances of banter among mock-fighting opponents, but also actual terms of endearment that can be used among peers following a seemingly subversive, yet ultimately conventionalized strategy. This has long been known also from other contexts. Upper middle-class men in white collar jobs running into old friends, whom they ostentatiously refer to as ``(you) son of a bitch''. Probably not all offensive vocabulary qualifies for use in such situations. The word ``son of a bitch'' does, which indicates a certain degree of conventionalization despite the fact that the two men would probably assume that their choice of word is an individual thing among them as intimates. The use of a term like ``son of a bitch'' in such contexts requires the presence of others, a third party, in order to work properly. It is questionable that the men would greet each other like that if they met while hiking in a forest with no other people around. They are part of a performed demonstration of one's social links, positioning oneself with regard to the other. 

Such terms are functionally ambivalent then. Their use as terms of endearment requires a canvas, a place for the projection of playful behavior. A stage. They are performative emblems in a fairly literal sense. This potential for ambivalence underscores the need to move away from dyadically oriented research. Speakers act not only according to their immediate interlocutor's role in communication; they monitor themselves and factor in ingroup practices, expectations, norms, and aesthetic preferences.

We assume, but cannot be entirely sure, that a performative motivation also underlies what happened to Hilal at his workplace one day. A regular customer shared his negative views on foreigners in Germany. That customer's xenophobic statement right in front of him as a foreigner, actually addressing him directly in the conversation, is hard to interpret for Hilal. The customer had always appeared to like him, they had often chatted, they had known each other for a long time. Does this mean, the customer exempts him from the group of foreigners about whose presence he complains? Difficult to imagine because the customer knows Hilal so well, and he can hardly have overlooked the fact that Hilal is a young man who grew up in \isi{Morocco}. So, is it intended to be provocative? The customer leads Hilal into a discursive double-bind: if Hilal ignored the statement, he'd be complicit; if he started arguing, he'd step into the trap of accepting being associated with that group as a meaningful categorization schema. 

Marcus' account of being discriminated against in the student cafeteria failed to spark much sympathy among the other conversation participants. This may very well be an effect of him being the weakest participant in terms of conversational space, impact, and attention that others would accord to him during the recording session. Marcus started the conversation resolutely, telling us about three young men who were trying to jump the queue and called him ``Alman'' -- a term that he experienced as offensive: 
\begin{quote}
\textit{Drei Türkischsprachige haben mich als Alman beleidigt, weil, das war halt in der Mensa und die haben sich halt vorgedrängelt, und ich hab halt gemeint, sie sollen da hinten sich anstellen, weil die Mensa halt voll war. Und dann haben die halt gesagt, ja, Scheiß Alman immer so die Regeln beachten, so … Deswegen …}

`Three \ili{Turkish}-speaking guys offended me calling me Alman, because, well, that was in the student cafeteria, they jumped the queue, so I am like get into the queue at the back, because the cafeteria was full. And then they are like, yes, fucking Alman always following the rules, like that … that's why…' (Marcus [00:01:28--00:01:44]) 
\end{quote}

Hilal and Lamin looked at Marcus skeptically. Apparently, they disagreed with Marcus on whether to interpret this as discriminatory behavior. Without explicitly saying anything, they conveyed their impression that this may have been insignificant banter. Marcus relativizes his initially strong statement imbued with genuine anger that one could tell from the tone of his voice. In a loud, but more conciliatory voice, he says: 
\begin{quote}
\textit{Ich mein, ok, vielleicht ist es so'n Alman-Ding, dass man so sich halt hintenanstellt, und nicht nur Ausländer und Migranten drängeln sich vor, Deutsche drängeln sich auch vor, ich kenn auch \textnormal{[\textit{zu sagen}] (unclear)} er soll nach hinten gehen und gleich beleidigt werden … iss schon komisch, eigentlich kennt man das nicht.}

`I mean, ok, maybe it is an Alman thing that one respects the queue, and not only foreigners and migrants jump the queue; I also know [unclear] but … getting insulted for telling someone to join the queue … it is weird, usually one doesn't experience such things.' (Marcus [00:01:52])
\end{quote}

Hilal seizes the opportunity to tell us about his own experiences. He steers the conversation in the direction of a discussion of lexical meaning and etymologies. With regard to what Marcus experienced as an offense, Lamin cuts in: \textit{Aber er ist Deutscher!} `But he \textit{is} \ili{German}!' (Lamin [00:10:19]), implying that the term is actually neutral in that it simply means `\ili{German}' in \ili{Turkish}. By insisting that the term is semantically adequate, he either fails to understand or is unwilling to concede that Marcus's negative experience is primarily about feeling ``pragmatically cornered'' in the cafeteria situation: After Marcus had asked the \ili{Turkish}-speaking men to join the queue, they gave in. Seemingly, Marcus had imposed his will. But the three men managed to make the whole situation about something else than respecting the queue. How was Marcus supposed to react in that situation? If he were to complain about being called an Alman, he would underscore their point -- complaining about being called an Alman is very Alman-like behavior. If he were to abstain from reacting to their verbal offense, they would have effectively silenced him. 

In the group conversation, we are dealing partly with a kind of common misunderstanding. The speaker's motives are not understood or interpreted properly. Marcus's use of modal particles is significant in this regard. He frequently adds \textit{halt} to his sentences, which signals that he is aware of possible objections to his ``claim of being discriminated against'', but appeals to the listeners to show sympathy for his position. The modal value of \textit{halt} is complex. There are various layers of meaning to it, and its pragmatic function is not comprehensively captured by this short description. However, it suffices in the current context to make the central point, namely that Hilal and Lamin failed to understand Marcus's intentions. It is not simply a misunderstanding but involves an instance of failed accommodation. Marcus's very frequent use of \textit{halt} is also intended to make his way of speaking more accessible, less formal, and closer to an encounter outside the university. The interpretation of modal particles in informal registers is fairly complicated, though, and fails the accommodative strategy he intended to pursue.

This happens against the backdrop of two possible explanations concerning especially Hilal's, but also Lamin's reaction. Either they did not understand (then there is no reason to be upset), or their failure to react empathically in a more accommodative way is interpreted as intentional nonaccommodation. Ultimately, Marcus's explicit attempt at fixing the communicative misunderstanding fails. The topic is aborted; Hilal and Lamin take the stage and turn to their own discriminatory experiences. Marcus takes on a more passive role, a little annoyed and a little disappointed. Later, he remains fairly silent, cautious, and slightly marginalized in the conversation. Being a native speaker of \ili{German} might have afforded him a communicative advantage. Yet, the framing of the conversation in terms of experiences of discrimination made the situation difficult for him, given his status as ``non-migrant''. His silence does not figure in the transcriptions. It is, therefore, hard to pinpoint where his reservation is motivated by positive intent to listen to others, politeness, or withdrawal from the thematic thrust of the conversation. At that level, the fact that two of the authors were present as interlocutors but also as critical observers enriches interpretations that would have been more restrained in the case of an analysis based solely on a narrow transcription of the recording. We believe that what Hilal and Marcus experienced relates to the most powerful mechanism behind ethnophaulisms. Their use, together with the strong meaning they carry, came into being in particular sociohistorical contexts. They are obviously offensive, serve as powerful insults and can affect individuals both collectively and on a personal level. Yet, even in our very cursory treatment of explicit derogatory labels, we were able to show that matters are more complex than that. Ethnophaulisms build on what is regarded as the Other's negative attributes, but rather than explaining the Other, they reflect an ingroup desire to differentiate oneself from the outgroup. But they can be subverted, either when individuals adopt them as self-designations among peers or when used in other performative functions, for instance, in staged performances or in conversations across group boundaries when their offensive denotations are canceled by the appropriate social conditions. In fact, their use can enhance social ties because metapragmatically, their potentiality to be used in such contexts signals robust social bonds -- running a relatively great risk of miscommunication, though. Jokes may easily fail, humorous provocations can cause misapprehension, and ingroup use of such terms possibly still spark negative feelings among addressees -- a fine line.

\section{Closing the conversation: nonaccommodation and critical observations on agency}\label{CH1.1.ss.5}

With our chapter, we pursued three main interests. A central descriptive aim was to document an open conversation on experiences of being discriminated against, as narrated primarily by male migrants from North Africa. Secondly, we wanted to push our methodological approach into a more multilateral conversation in a group. We were hoping for insights to emerge from the fact that the participants would react to each other in more varied ways, given that their constellation was mixed in terms of social identities, cultural and linguistic backgrounds and roles in the research process. Finally, our contribution was motivated by a theoretical concern: our interest in contemplating how CAT and critical frameworks relate to one another, hoping for mutual benefit. Where do we stand with regard to these three aims? In this closing discussion, we shall try to address this question in reverse order -- theory, method, data -- because we would like to end the conversation by evaluating the voices of those who shared their experiences with us.

A close reading of the group interview that forms the basis of this chapter and its fairly narrow transcription prepared the grounds for a description in terms of the central concepts of CAT. The different episodes that we zoomed in on provide empirical support for CAT. As a heuristic tool, framing the various situations that were narrated by the three young men in terms of intergroup scenarios afforded analytical dimensions that proved useful, such as paying attention to the granularity of identity perceptions (individual or social?), or distinguishing between ingroup versus outgroup interaction. In the narrated events, these categories turned out to be meaningful and could be accessed, although the original communicative situations portrayed in the narrations were not directly observed by us as researchers nor by the other participants in the conversation. This is not necessarily a drawback. In fact, the discussions that followed the narration of individual episodes highlighted the most significant properties relying on ideas and notions that lend themselves to an understanding in terms of CAT, e.g., when Hilal and Lamin disagreed on how to assess particular roles and motives in the police episode. This underscores how intuitive a great number of the proposals in CAT as a conceptual framework are, including the fact that it takes into consideration subjective dimensions. People react to communicative behaviors as they perceive them, which is not necessarily identical to observable communicative behaviors. What's more, perceptions and interpretations of communicative behaviors may differ greatly among interlocutors.

The police episode indicates that how we evaluate social interactions and our conversational partners depends not only on our idiosyncratic characteristics but largely on our group membership and the relevant group to which we believe our conversational partner belongs. This is an important point -- alignment does not happen automatically but depends on speakers' expectations, world knowledge and assumptions concerning one's interlocutors. In that sense, group membership is as much presupposed as it is contested, denied or renegotiated in communicative events. 

Communicative needs and desires may call for speech adjustments (e.g., a native speaker adjusting their speech so that they are understood by a non-native speaker; using the right code, e.g., standard variety vs. regional \isi{dialect}, high/prestige language vs. low language), but they also involve affective needs. Acknowledging affective needs may involve something as simple as responding to a greeting. The policeman's unwillingness to adjust and shift to Hilal's preferred language is a clear violation of Hilal's affective expectation and need, based on what, to him, is a fundamental issue in terms of social identity and \isi{belonging}, and therefore sparks strong emotional reactions. 

During communication, people negotiate personal and social identities. Conversational partners or even ``silent speakers'' (see the episode in \sectref{CH1.1.ss.4.2}, the bag in the supermarket) often categorize the Other based on stereotypes rather than on actual qualities of an individual -- their personal identity -- and/or of their group, in other words, their social identity. When that happens, the communicative event is perceived as nonaccommodative, as it does not meet the needs and desires of the recipient; in case of particularly confrontational behaviors (involving, for example, discrimination and hostile communicative practices), as counteraccommodative. It is important to stress that this kind of counteraccommodation does not require the person who causes it to act consciously: Counteraccommodation is not the same as an intentional offense, as illustrated by the pain caused by specific nonverbal reactions of someone to another person's mere presence, based on harmful stereotypes. They are psychologically dissociative. Of the three kinds of interactions that this chapter illustrated, the use of abusive terms and ethnophaulisms appears as the clearest case of such psychological dissociation. Shouting at someone ``Go back to your country!'' is certainly intentional, hostile, and counteraccommodative. Yet, the use of derogatory language among peers can work differently. Slightly surprising, albeit not uncommon, are instances of allegedly derogatory expressions that do not work as an offense in a particular environment and conversational encounter. 

Especially with regard to such apparent dissonances, we believe that a critical analysis of the speech event that we recorded can add significantly to our understanding of additional layers of meaning. The context of a group interview has proved helpful, because any dissonances, incidents of failed communication and nonaccommodation would not simply be reported by an individual interviewee, but also discussed at a meta-level in the group. The relative heterogeneity of the group did not only add different viewpoints as in a panel discussion. It helped a long way in making the conversation natural and meaningful to all participants. Of course, such a format is not without its challenges. Often, answers were provided that did not respond to the questions asked. Conversational maxims were not always followed. In this regard, our group discussion resembled everyday communication perhaps more than an interview format. 

Our decision to form a group of participants who would share some, but differ in terms of other attributes that are significant in terms of identity construal meant that the participants had to increase their efforts in figuring out who would understand what in terms of their respective commonalities and differences. The constellation enabled a conversation of shared responsibilities and a discussion of issues that are relevant to the participants at a meta-level. At the same time, it brought to the fore the emotional impact that the reported incidents of behavior that was perceived as discriminatory had on the participants. 

In social interactions, people have expectations, desires and intentions. These influence the way they communicate and the way they interpret communication. All the interactions we had were intergroup, implying a constant game of approval and rejection. Being discriminated against, or rather having the perception that that is the case, has negative consequences not only because social distances are reinforced but also because receivers are called to face situations and behaviors that lead to emotional burden and distress. 

Contrary to what one might expect, dealing with the anger that can be caused by nonaccommodation does not necessarily work along an intensity scale from overaccommodation to underaccommodation, all the way to outright counteraccommodation. Even an outright insult, despite its negative intent and the discomfort it causes the addressee, may be easier to handle than other forms of nonaccommodation. The contributions to our conversation by Hilal, Lamin and Marcus all coincide in that they show quite clearly what is most vexing about nonaccommodation and its interpretation as discriminatory behavior: a sense of powerlessness, lack of choices of action, being effectively silenced and thereby having one's agency curtailed. 

The context of the group conversation was artificial. Not all of the participants had known one another before. The recording session took place in a somewhat sterile meeting room on the university premises. The participants found themselves in a situation that imposed a generally calm, controlled and respectful interaction. Still, subtle power play in terms of how discursive positions were built during the conversation could be observed. Interestingly, the power asymmetries did not follow the obvious cultural or linguistic boundaries represented by the participants.

Overall, noticeable nonaccommodation did not occur frequently during the group conversation that we recorded and analyzed, probably because of the contrived context of data recording in a university setting. So, one may wonder whether the approach of recording speech events in such ``semi-natural'' conditions is fruitful then. We believe it is because of certain unexpected observations -- e.g., nonaccommodation among (arguably) ingroup members according to different levels of proficiency\is{Language!proficiency}. 

CAT proved particularly helpful with regard to those portions of the conversation characterized largely by dialogic interaction between Hilal and Lamin. By then, we had recorded the conversation for about fifteen minutes. Their conversation (mostly a dialogue at this stage) seemed genuine, with interruptions, contradictions, and opposition to what the interlocutor said. It was respectful, though. After Hilal ended the rather monologic narration of his experience in \isi{Morocco}, the conversational chunks remained fairly long; turn-taking did not accelerate too much. Lamin remains fairly calm but talks back to Hilal, speaking in a more controlled (and possibly therefore closer to standard) way in \ili{German}. Hilal is angered and trapped by Lamin's skillful use of a question strategy in order to make his point about the policemen (possibly speaking \ili{Darija} rather than Tarifit\il{Berber!Tarifit} for no particular reason): \textit{Die sind offener […], können halt beides, Arabisch und Berbisch} `They are more open … they simply know both, \ili{Arabic} and \ili{Berber}' (Lamin [00:17:50]). Hilal insists it is because of the rules that the policeman mentioned and asks Lamin provocatively: \textit{Weißt du, was Regeln sind?} `Do you know what rules are?' (Hilal [00:18:14]). But Lamin turns the strategy around, using a sequence of rhetorical questions himself: 
\begin{quote}
\textit{Warte, er ist \ili{Berber}, er redet mit seinem Kollegen Berbisch\il{Berbisch|see{Berber}}, Und dann provoziert er einen \ili{Berber}? Wegen Regeln von seinem Chef -- der ist Araber oder \ili{Berber}?}

`Wait, he is \ili{Berber}, he uses \ili{Berber} with his colleague, and then he goes on to provoke a \ili{Berber}? Because of rules laid out by his boss? Who is Arab or Berber?' (Lamin [00:18:37--00:18:52]) 
\end{quote}

Hilal claims the boss is definitely Arab, but Lamin counters, \textit{Also man weiß das im Vorhinein?} `Oh, one knows that in advance?' (Lamin [18:52]), discrediting Hilal's account forcefully with his rhetorical questions.

At that stage, the conversation is characterized by playful one-upmanship. It is extremely hard to tell whether Lamin's conversational advantage is linked to a possibly better command of \ili{German}, to a situational advantage, or to greater communicative skill, which may but does not have to be tied to different degrees of proficiency\is{Language!proficiency} in \ili{German} as a language that both of them learned only as adults. Hilal is emotionally more involved because he talks about a personal memory that caused him pain and because he is arguably more sensitive to discrimination on the basis of being \ili{Berber}, attested to by his statements and fervor throughout the conversation. In contrast, Lamin remains calm and more detached, which gives him an air of superiority at times. The communicative misalignments that we had expected to find across the Tamazight\il{Berber!Tamazight}/\ili{German} divide occur at a different level between Hilal and Lamin. Both of them apply strategies that are not uncommon in conversations among (especially younger) North African men -- being provocative, defying one's interlocutor, challenging their views in a provocative way, shielding oneself behind apparently innocent questions and sometimes explicit use of excusatory formulae. 

Such communicative situations require a dual perspective. Like in this case, it is often hard to tell when we are dealing with (possibly unintended, yet inappropriate) nonaccommodation and when the interlocutors' statements break conversational maxims of collaboration on purpose. In our case of a contrived conversation in a university setting, certain limitations may have affected the richness of the conversation texture. We hope to have been able to show nevertheless that there is added value in adopting a dual perspective, drawing on CDA and CAT. It is hard to pinpoint where one of the frameworks would work better and prove more beneficial to the interpretation. Given the broad overlap in terms of their general tenets, we believe that remaining somewhat eclectic in this regard adds to the interpretation by broadening our toolset and enhancing conceptual awareness in what is ultimately an ethnographic approach resting on two theoretical pillars. 

\section*{Abbreviations}
\begin{tabularx}{.75\textwidth}{@{} l l @{}}
CAT&\isi{Communication Accommodation Theory}\\
CDA&\isi{Critical Discourse Analysis\is{Critical Discourse Analysis (CDA)}}\\
\end{tabularx}

% \section*{Acknowledgements}
% \citet{Nordhoff2018} is useful for compiling bibliographies
%\section*{Contributions}
%John Doe contributed to conceptualization, methodology, and validation. 
%Jane Doe contributed to writing of the original draft, review, and editing.

\sloppy
\printbibliography[heading=subbibliography,notkeyword=this]
\end{document}
